\documentclass[11pt]{article}
\usepackage[margin=1in]{geometry}
\usepackage{booktabs}
\usepackage{array}
\usepackage{longtable}
\usepackage[T1]{fontenc}
\usepackage[utf8]{inputenc}
\title{QR-HybridFin Methodology Appendix}
\author{}
\date{}
\begin{document}
\maketitle
\section*{Overview}
This appendix condenses the transaction-level simulator into paper-style tables and captions.
It is derived from the layered model implemented in \texttt{poc/qr\_hybridfin\_methodology\_sim.py}.
\section*{Table A1. Core Simulation Outcomes}
\begin{table}[h]
\centering
\begin{tabular}{@{}p{0.62\linewidth}r@{}}
\toprule
Metric & Value \\
\midrule
Adaptive trace mean latency & 106.258 ms \\
Adaptive trace p95 latency & 285.982 ms \\
Adaptive classical selections & 1 \\
Adaptive static-hybrid selections & 5587 \\
Adaptive ML-optimized selections & 4412 \\
Best throughput at 50 nodes & 1012.64 TPS \\
Worst cross-border mean settlement & 242.30 ms \\
Classical ablation mean latency & 105.506 ms \\
Static-hybrid ablation mean latency & 111.582 ms \\
ML-optimized ablation mean latency & 108.157 ms \\
\bottomrule
\end{tabular}
\end{table}
\section*{Table A2. Derived Interpretation}
\begin{table}[h]
\centering
\begin{tabular}{@{}p{0.34\linewidth}p{0.58\linewidth}@{}}
\toprule
Observation & Interpretation \\
\midrule
Classical fallback remains available & Low-risk domestic transactions can stay on the legacy path during migration. \\
Static hybrid dominates mid-load traffic & The agility manager shifts to the hybrid mode as load rises. \\
ML-optimized covers cross-border/high-threat flows & The policy layer pushes the stronger profile where risk or WAN friction is highest. \\
Throughput remains in the paper's target range & The batch service model preserves the 1180 / 965 / 1015 TPS calibration envelope. \\
\bottomrule
\end{tabular}
\end{table}
\section*{Figure Captions}
\begin{itemize}
\item Figure A1. QR-HybridFin architecture diagram showing the cryptographic, blockchain, ML agility, and network layers connected by the transaction flow.
\item Figure A2. Cryptographic latency comparison anchored to the paper's Table I values and ML-optimized overhead reduction.
\item Figure A3. Wholesale throughput comparison showing the calibrated classical, static-hybrid, and ML-optimized TPS points.
\item Figure A4. Throughput scaling across 50-500 nodes under the paper's CBDC deployment scenario.
\item Figure A5. Latency versus payload size generated from the layered transaction cost model.
\item Figure A6. Cross-border settlement latency under Mininet + tc WAN emulation with 100-300 ms links.
\item Figure A7. Quantum resistance comparison for classical and hybrid schemes under the paper's Shor/Grover discussion.
\end{itemize}
\section*{Methodological Notes}
\begin{itemize}
\item The per-transaction breakdown records cryptographic, blockchain, routing, queueing, and ML-agility terms separately.
\item The adaptive selector uses transaction threat level, current load, and cross-border status to choose the cryptographic profile.
\item Cross-border and ablation summaries use the same cost model, so the comparison is generated from one simulation pipeline rather than a replot of static values.
\end{itemize}
\end{document}
